% python classes slides - classes_introduction
% (c) 2014 Kostiantyn Danylov aka koder 
% koder.mail@gmail.com
% distributed under CC-BY licence
% http://creativecommons.org/licenses/by/3.0/deed.en

\documentclass{article}
% XeLaTeX
\usepackage{xltxtra}
\usepackage{xunicode}
\usepackage{listings}
\usepackage[landscape]{geometry}

% Fonts
\setmainfont{DejaVu Sans} %{Arial}
\newfontfamily\cyrillicfont{Nimbus Roman No9 L} %{Arial}
\setmonofont{Courier New}
%\setmonofont{Ubuntu Mono}

%\setmonofont{DejaVu Sans Mono}

% Lang
\usepackage{polyglossia}
\setmainlanguage{russian}
\setotherlanguage{english}
\usepackage[dvipsnames,table]{xcolor}


\ifx\pdfoutput\undefined
\usepackage{graphicx}
\else
\usepackage[pdftex]{graphicx}
\fi

\lstset{
	language=python,
	keywordstyle=\color{Emerald},%\texttt, 
	commentstyle=\color{OliveGreen},%\texttt,
	stringstyle=\color{Bittersweet},%\texttt,
	tabsize=4,
	numbers=left,
	xleftmargin=10pt,
	morekeywords={with,as},	
	numberstyle=\large,
	%identifierstyle=\texttt,
	%basicstyle=\texttt,
}

\usepackage{hyperref}

\hypersetup{
	colorlinks=true,
	urlcolor=blue
}

\usepackage{float}
%\floatstyle{boxed} 
%\restylefloat{figure}
\usepackage[normalem]{ulem}


\makeatletter
\def\PY@reset{\let\PY@it=\relax \let\PY@bf=\relax%
    \let\PY@ul=\relax \let\PY@tc=\relax%
    \let\PY@bc=\relax \let\PY@ff=\relax}
\def\PY@tok#1{\csname PY@tok@#1\endcsname}
\def\PY@toks#1+{\ifx\relax#1\empty\else%
    \PY@tok{#1}\expandafter\PY@toks\fi}
\def\PY@do#1{\PY@bc{\PY@tc{\PY@ul{%
    \PY@it{\PY@bf{\PY@ff{#1}}}}}}}
\def\PY#1#2{\PY@reset\PY@toks#1+\relax+\PY@do{#2}}

\expandafter\def\csname PY@tok@gd\endcsname{\def\PY@tc##1{\textcolor[rgb]{0.63,0.00,0.00}{##1}}}
\expandafter\def\csname PY@tok@gu\endcsname{\let\PY@bf=\textbf\def\PY@tc##1{\textcolor[rgb]{0.50,0.00,0.50}{##1}}}
\expandafter\def\csname PY@tok@gt\endcsname{\def\PY@tc##1{\textcolor[rgb]{0.00,0.25,0.82}{##1}}}
\expandafter\def\csname PY@tok@gs\endcsname{\let\PY@bf=\textbf}
\expandafter\def\csname PY@tok@gr\endcsname{\def\PY@tc##1{\textcolor[rgb]{1.00,0.00,0.00}{##1}}}
\expandafter\def\csname PY@tok@cm\endcsname{\let\PY@it=\textit\def\PY@tc##1{\textcolor[rgb]{0.25,0.50,0.50}{##1}}}
\expandafter\def\csname PY@tok@vg\endcsname{\def\PY@tc##1{\textcolor[rgb]{0.10,0.09,0.49}{##1}}}
\expandafter\def\csname PY@tok@m\endcsname{\def\PY@tc##1{\textcolor[rgb]{0.40,0.40,0.40}{##1}}}
\expandafter\def\csname PY@tok@mh\endcsname{\def\PY@tc##1{\textcolor[rgb]{0.40,0.40,0.40}{##1}}}
\expandafter\def\csname PY@tok@go\endcsname{\def\PY@tc##1{\textcolor[rgb]{0.50,0.50,0.50}{##1}}}
\expandafter\def\csname PY@tok@ge\endcsname{\let\PY@it=\textit}
\expandafter\def\csname PY@tok@vc\endcsname{\def\PY@tc##1{\textcolor[rgb]{0.10,0.09,0.49}{##1}}}
\expandafter\def\csname PY@tok@il\endcsname{\def\PY@tc##1{\textcolor[rgb]{0.40,0.40,0.40}{##1}}}
\expandafter\def\csname PY@tok@cs\endcsname{\let\PY@it=\textit\def\PY@tc##1{\textcolor[rgb]{0.25,0.50,0.50}{##1}}}
\expandafter\def\csname PY@tok@cp\endcsname{\def\PY@tc##1{\textcolor[rgb]{0.74,0.48,0.00}{##1}}}
\expandafter\def\csname PY@tok@gi\endcsname{\def\PY@tc##1{\textcolor[rgb]{0.00,0.63,0.00}{##1}}}
\expandafter\def\csname PY@tok@gh\endcsname{\let\PY@bf=\textbf\def\PY@tc##1{\textcolor[rgb]{0.00,0.00,0.50}{##1}}}
\expandafter\def\csname PY@tok@ni\endcsname{\let\PY@bf=\textbf\def\PY@tc##1{\textcolor[rgb]{0.60,0.60,0.60}{##1}}}
\expandafter\def\csname PY@tok@nl\endcsname{\def\PY@tc##1{\textcolor[rgb]{0.63,0.63,0.00}{##1}}}
\expandafter\def\csname PY@tok@nn\endcsname{\let\PY@bf=\textbf\def\PY@tc##1{\textcolor[rgb]{0.00,0.00,1.00}{##1}}}
\expandafter\def\csname PY@tok@no\endcsname{\def\PY@tc##1{\textcolor[rgb]{0.53,0.00,0.00}{##1}}}
\expandafter\def\csname PY@tok@na\endcsname{\def\PY@tc##1{\textcolor[rgb]{0.49,0.56,0.16}{##1}}}
\expandafter\def\csname PY@tok@nb\endcsname{\def\PY@tc##1{\textcolor[rgb]{0.00,0.50,0.00}{##1}}}
\expandafter\def\csname PY@tok@nc\endcsname{\let\PY@bf=\textbf\def\PY@tc##1{\textcolor[rgb]{0.00,0.00,1.00}{##1}}}
\expandafter\def\csname PY@tok@nd\endcsname{\def\PY@tc##1{\textcolor[rgb]{0.67,0.13,1.00}{##1}}}
\expandafter\def\csname PY@tok@ne\endcsname{\let\PY@bf=\textbf\def\PY@tc##1{\textcolor[rgb]{0.82,0.25,0.23}{##1}}}
\expandafter\def\csname PY@tok@nf\endcsname{\def\PY@tc##1{\textcolor[rgb]{0.00,0.00,1.00}{##1}}}
\expandafter\def\csname PY@tok@si\endcsname{\let\PY@bf=\textbf\def\PY@tc##1{\textcolor[rgb]{0.73,0.40,0.53}{##1}}}
\expandafter\def\csname PY@tok@s2\endcsname{\def\PY@tc##1{\textcolor[rgb]{0.73,0.13,0.13}{##1}}}
\expandafter\def\csname PY@tok@vi\endcsname{\def\PY@tc##1{\textcolor[rgb]{0.10,0.09,0.49}{##1}}}
\expandafter\def\csname PY@tok@nt\endcsname{\let\PY@bf=\textbf\def\PY@tc##1{\textcolor[rgb]{0.00,0.50,0.00}{##1}}}
\expandafter\def\csname PY@tok@nv\endcsname{\def\PY@tc##1{\textcolor[rgb]{0.10,0.09,0.49}{##1}}}
\expandafter\def\csname PY@tok@s1\endcsname{\def\PY@tc##1{\textcolor[rgb]{0.73,0.13,0.13}{##1}}}
\expandafter\def\csname PY@tok@sh\endcsname{\def\PY@tc##1{\textcolor[rgb]{0.73,0.13,0.13}{##1}}}
\expandafter\def\csname PY@tok@sc\endcsname{\def\PY@tc##1{\textcolor[rgb]{0.73,0.13,0.13}{##1}}}
\expandafter\def\csname PY@tok@sx\endcsname{\def\PY@tc##1{\textcolor[rgb]{0.00,0.50,0.00}{##1}}}
\expandafter\def\csname PY@tok@bp\endcsname{\def\PY@tc##1{\textcolor[rgb]{0.00,0.50,0.00}{##1}}}
\expandafter\def\csname PY@tok@c1\endcsname{\let\PY@it=\textit\def\PY@tc##1{\textcolor[rgb]{0.25,0.50,0.50}{##1}}}
\expandafter\def\csname PY@tok@kc\endcsname{\let\PY@bf=\textbf\def\PY@tc##1{\textcolor[rgb]{0.00,0.50,0.00}{##1}}}
\expandafter\def\csname PY@tok@c\endcsname{\let\PY@it=\textit\def\PY@tc##1{\textcolor[rgb]{0.25,0.50,0.50}{##1}}}
\expandafter\def\csname PY@tok@mf\endcsname{\def\PY@tc##1{\textcolor[rgb]{0.40,0.40,0.40}{##1}}}
\expandafter\def\csname PY@tok@err\endcsname{\def\PY@bc##1{\setlength{\fboxsep}{0pt}\fcolorbox[rgb]{1.00,0.00,0.00}{1,1,1}{\strut ##1}}}
\expandafter\def\csname PY@tok@kd\endcsname{\let\PY@bf=\textbf\def\PY@tc##1{\textcolor[rgb]{0.00,0.50,0.00}{##1}}}
\expandafter\def\csname PY@tok@ss\endcsname{\def\PY@tc##1{\textcolor[rgb]{0.10,0.09,0.49}{##1}}}
\expandafter\def\csname PY@tok@sr\endcsname{\def\PY@tc##1{\textcolor[rgb]{0.73,0.40,0.53}{##1}}}
\expandafter\def\csname PY@tok@mo\endcsname{\def\PY@tc##1{\textcolor[rgb]{0.40,0.40,0.40}{##1}}}
\expandafter\def\csname PY@tok@kn\endcsname{\let\PY@bf=\textbf\def\PY@tc##1{\textcolor[rgb]{0.00,0.50,0.00}{##1}}}
\expandafter\def\csname PY@tok@mi\endcsname{\def\PY@tc##1{\textcolor[rgb]{0.40,0.40,0.40}{##1}}}
\expandafter\def\csname PY@tok@gp\endcsname{\let\PY@bf=\textbf\def\PY@tc##1{\textcolor[rgb]{0.00,0.00,0.50}{##1}}}
\expandafter\def\csname PY@tok@o\endcsname{\def\PY@tc##1{\textcolor[rgb]{0.40,0.40,0.40}{##1}}}
\expandafter\def\csname PY@tok@kr\endcsname{\let\PY@bf=\textbf\def\PY@tc##1{\textcolor[rgb]{0.00,0.50,0.00}{##1}}}
\expandafter\def\csname PY@tok@s\endcsname{\def\PY@tc##1{\textcolor[rgb]{0.73,0.13,0.13}{##1}}}
\expandafter\def\csname PY@tok@kp\endcsname{\def\PY@tc##1{\textcolor[rgb]{0.00,0.50,0.00}{##1}}}
\expandafter\def\csname PY@tok@w\endcsname{\def\PY@tc##1{\textcolor[rgb]{0.73,0.73,0.73}{##1}}}
\expandafter\def\csname PY@tok@kt\endcsname{\def\PY@tc##1{\textcolor[rgb]{0.69,0.00,0.25}{##1}}}
\expandafter\def\csname PY@tok@ow\endcsname{\let\PY@bf=\textbf\def\PY@tc##1{\textcolor[rgb]{0.67,0.13,1.00}{##1}}}
\expandafter\def\csname PY@tok@sb\endcsname{\def\PY@tc##1{\textcolor[rgb]{0.73,0.13,0.13}{##1}}}
\expandafter\def\csname PY@tok@k\endcsname{\let\PY@bf=\textbf\def\PY@tc##1{\textcolor[rgb]{0.00,0.50,0.00}{##1}}}
\expandafter\def\csname PY@tok@se\endcsname{\let\PY@bf=\textbf\def\PY@tc##1{\textcolor[rgb]{0.73,0.40,0.13}{##1}}}
\expandafter\def\csname PY@tok@sd\endcsname{\let\PY@it=\textit\def\PY@tc##1{\textcolor[rgb]{0.73,0.13,0.13}{##1}}}

\def\PYZbs{\char`\\}
\def\PYZus{\char`\_}
\def\PYZob{\char`\{}
\def\PYZcb{\char`\}}
\def\PYZca{\char`\^}
\def\PYZam{\char`\&}
\def\PYZlt{\char`\<}
\def\PYZgt{\char`\>}
\def\PYZsh{\char`\#}
\def\PYZpc{\char`\%}
\def\PYZdl{\char`\$}
\def\PYZti{\char`\~}
% for compatibility with earlier versions
\def\PYZat{@}
\def\PYZlb{[}
\def\PYZrb{]}
\makeatother


\begin{document}
\LARGE
%-------------------------------------------------------------------------------
\begin{center} ООП \end{center}
    \begin{itemize}
         \item Повторное использование кода
         \item Инкапсуляция
         \item Наследование
         \item Полиморфизм
     \end{itemize} 
\newpage
%-------------------------------------------------------------------------------
\begin{center} \includegraphics[scale=5]{images/Hamatsa_shaman2.jpg} \end{center}
\newpage

%-------------------------------------------------------------------------------
\begin{center} Подсчет среднеарифметического \end{center}
\begin{lstlisting}
    def mean_int(numbers):
        res = 0
        for cnum in numbers:
            res += cnum
        return res / len(numbers)
\end{lstlisting}
\newpage

%-------------------------------------------------------------------------------
\begin{center} Нужно добавить поддержку рациональных чисел \end{center}
Рациональное число - пара (числитель, знаменатель). Будем передавать
рациональные числа в виде кортежа.

$$
\frac{a}{b} + \frac{c}{d} = \frac{ad + bc}{bd}\,\,;
\quad
\frac{a}{b} / c = \frac{a}{cd}
$$
\newpage

%-------------------------------------------------------------------------------
\begin{center} Среднее рациональные \end{center}
\begin{lstlisting}
    # rational = (num, denom)
    def mean_rational(numbers):
        num, denom = (0, 1)
        for cnum, cdenom in numbers:
            num = num * cdenom + denom * cnum
            denom *= cdenom            
        return (num, denom * len(numbers))
\end{lstlisting}
\newpage

%-------------------------------------------------------------------------------
\begin{center} Среднее \end{center}
\begin{lstlisting}
    def mean(numbers):
        assert len(numbers) != 0
        if isinstance(numbers[0], (int, long, float)):
            return mean_int(numbers)
        else:
            return mean_rational(numbers)
\end{lstlisting}
\newpage

%-------------------------------------------------------------------------------
\begin{center} Добавляем матрицы \end{center}
\begin{lstlisting}
    def mean(numbers):
        assert len(numbers) != 0
        if isinstance(numbers[0], (int, long, float)):
            return mean_int(numbers)
        elif ismatrix(numbers[0]):
            return mean_matrix(numbers)
        else:
            return mean_rational(numbers)
\end{lstlisting}
\newpage

%-------------------------------------------------------------------------------
\begin{center} "Идея" СА \end{center}
\begin{lstlisting}
    def mean_int(numbers):
        res = 0
        for cnum in numbers:
            res += cnum
        return res / len(numbers)
\end{lstlisting}
\newpage

%-------------------------------------------------------------------------------
\begin{center}  СА \end{center}
\begin{lstlisting}
    add_num = lambda x, y: x + y
    add_rational = lambda (x1, y1), (x2, y2): \
                    (x1 * y2 + y1 * x2, y1 * y2)
    div_num = lambda x, y: x / y
    div_rational = lambda (x, y), z: (x, y * z)

    def mean(numbers, add_func, div_func):
        res = numbers[0]
        for cnum in numbers[1:]:
            res = add_func(res, cnum)
        return div_func(res, len(numbers))

    mean([1, 2, 3] , add_num, div_num)
\end{lstlisting}
\newpage

%-------------------------------------------------------------------------------
\begin{lstlisting}
    def mean(numbers, add_func, div_func):
        it = iter(numbers)
        res = next(it)
        for cnum in it:
            #...
\end{lstlisting}
\newpage

%-------------------------------------------------------------------------------
\begin{lstlisting}
    val = {'num':1, 
           'denom':2, 
           '__add__':add_rational,
           '__div__':div_rational}

    def add_rational(x, y):
        return  (x['num'] * y['denom'] + \
                 x['denom'] * y['num'], 
                 x['denom'] * y['denom'])

    def div_rational(x, y):
        return (x['num'], x['denom'] * y)
\end{lstlisting}
\newpage

%-------------------------------------------------------------------------------
\begin{lstlisting}
    def mk_rational(num, denom):
        return {'num':num, 
                'denom':denom, 
                'add':add_rational, 
                'div':div_rational}
\end{lstlisting}
\newpage

%-------------------------------------------------------------------------------
\begin{center} Нужно автоматически упрощать после операции \end{center}
\begin{lstlisting}
    def add_rational_auto_simpl(x, y):
        res = add_rational(x, y)
        nod = get_NOD(res['num'], res['denom'])
        res['num'] /= nod
        res['denom'] /= nod
        return res

    def mk_rational_auto_simpl(num, denom):
        return dict(num=num, 
                    denom=denom, 
                    __add__=add_rational_auto_simpl, 
                    __div__=div_rational_auto_simpl)
\end{lstlisting}
\newpage

%-------------------------------------------------------------------------------
\begin{center} На предыдущем слайде ошибка \end{center}
\newpage

%-------------------------------------------------------------------------------
\begin{lstlisting}
    def mk_rational(num, denom):
        return {'num':num, 
                'denom':denom, 
                '__class__':Rational}

    Rational = dict(__div__=div_rational,
                    __add__=add_rational,
                    __init__=mk_rational)
\end{lstlisting}
\newpage

%-------------------------------------------------------------------------------
\begin{center} Когда ООП \end{center}
\begin{itemize}
    \item Если есть участок кода, требующий определенного ограниченного 
          набора операций над входными данными
    \item Одновременно в программе могут быть несколько видов подходящих данных, 
          с различной функциональностью для реализации этого интерфейса
    \item Или группировки функций с общим глобальным состоянием
\end{itemize}
\newpage

%-------------------------------------------------------------------------------
\begin{center} Классы не предназначен для \end{center}
\begin{itemize}
    \item Группировки функций
    \item Группировки одной функции
    \item Если вы, ессно, используете нормальный ЯП
\end{itemize}
\newpage

%-------------------------------------------------------------------------------
\begin{center} Проблема \end{center}
\newpage

%-------------------------------------------------------------------------------
\begin{center} Проблема \end{center}
a + b для разных типов
\newpage

%-------------------------------------------------------------------------------
\begin{center} Среднее \end{center}
\begin{lstlisting}
    add = {(int, int):add_int_int, 
           (int, tuple):add_int_tuple,
           (tuple, int):add_tuple_int}

    def register_add(tp1, tp2, func):
        add[(tp1, tp2)] = func

    def mean(numbers):
        res = numbers[0]
        for cnum in numbers[1:]:
            res = add[(type(res), type(cnum))](res, cnum)
        return div[(type(res), int)](res, len(numbers)
\end{lstlisting}
\newpage

%-------------------------------------------------------------------------------
\end{document}
